% SPDX-License-Identifier: LPPL-1.3c
% Copyright (C) 2026 Christophe Jorssen
% Author and Current Maintainer: Christophe Jorssen <christophe.jorssen@gmail.com>
% LPPL maintenance status: maintained
% This file is part of luafemm. See LICENSE and LICENSES.md for its terms.
\section{A small alternating excitation on a DC bias}
\label{sec:biased-harmonic}

\subsection{Penetration of a small alternating field}
The static tangent in Section~\ref{sec:operating} relates a small current
change to a flux change. In a conducting iron rod, an oscillating perturbation
also depends on how far the alternating field penetrates.
The changing flux induces azimuthal eddy currents, which oppose penetration
and dissipate energy. Moreover, a strongly biased rod has a different tangent
permeability from a weakly biased one. Frequency and saturation must therefore
be considered together.

Keep the nonlinear \emph{DC} calculation and add a small sinusoidal
current. The package solves the complex linear response about that fixed DC
state. It does not repeatedly solve the nonlinear law along a finite hysteresis
loop. This distinction determines both the input and the interpretation of the
plots: all response fields exclude the DC field.

\subsection{Reuse the bias, choose a frequency}
The geometry and the capture mechanism are exactly those of the preceding
section. First declare the conducting material, coil and DC current,
then capture the converged solution at picture end.
In a second picture, select the captured name and add |frequency|:
\begin{codeexample}[code only]
\tikzset{ripple/.style 2 args={femm/problem={
  operating point=#1,permeability=incremental,frequency=#2,
  circuit currents={drive=1},circuit phases={drive=0}}}}
% The preceding DC picture has captured a point named rod-low.
\begin{tikzpicture}[ripple={rod-low}{200},
  femm/cache={file=rod-low-200,mode=auto}]
  \pic {femm geometry};
  \pic[femm/phase=0] {femm field};
  \pic[femm/currents={view=symbols,phase=0}] {femm eddy currents};
  \femmmeasure[name=terminal,quantity=circuit,circuit=drive,parameter=200]
  \femmmeasure[name=losses,parameter=200]
\end{tikzpicture}
\end{codeexample}
The |drive| circuit already belongs to the bias. Its current here is the
\emph{peak AC perturbation}, not the sum of DC and AC currents. A unit peak
current conveniently normalises the linear transfer functions. An actual
small peak amplitude $\epsilon$ multiplies the fields by $\epsilon$ and the
reported AC powers by $\epsilon^2$; it does not change impedance or inductance.
Unlisted circuits have zero AC current.

The two phase settings have different jobs. |circuit phases| changes the
physical source phasors and therefore the solve/cache signature.
|femm/phase|, or |phase| inside |femm/currents|, chooses the instant drawn from
an existing complex solution. It does not require another solve. All physical
phasors use the peak convention $\operatorname{Re}(\widehat u e^{i\omega t})$;
phases are in degrees.

\subsection{Two bias currents, one frequency sweep}
The complete example uses an axisymmetric rod of radius $5\,\mathrm{mm}$ and
length $24\,\mathrm{mm}$ inside a 100-turn stranded solenoid. The illustrated
half-section is rotated around $r=0$; its eddy currents close around that axis.
The conductivity is $2\,\mathrm{MS/m}$ and the B--H table is an illustrative
single-valued law, not a measured steel specification. The coil carries either
$5\,\mathrm A$ or $50\,\mathrm A$ DC. Each captured bias then supplies six
AC responses between 5 and $200\,\mathrm{Hz}$.

While arranging the drawing, try |\CoreMesh=.6| and only two frequencies.
Then use |.25|, followed by finer values, before interpreting small differences
in the curves. Refinement belongs to the \emph{bias} geometry: every response
inherits that exact mesh. A mesh sufficient for the DC field need not resolve
AC penetration. An approximate local skin depth is
$\delta=\sqrt{2/(\omega\mu_{\mathrm{diff}}\sigma)}$; in a nonuniform tensor
material it is only a guide. Check the smallest relevant penetration length
at the highest frequency and weakest differential reluctivity, as well as the
outer air boundary. Increasing contour counts or profile samples does not
refine the finite elements.

\luafemmexamplebegin{biased-harmonic}

The geometry below is shared by both bias currents. Each frequency response reuses its own bias mesh.

% manual example: biased-harmonic/1

\begin{codeexample}[setup code]
\definecolor{weakbias}{RGB}{20,85,135}
\definecolor{strongbias}{RGB}{185,65,45}
\newbox\BiasCalculation
\def\CoreMesh{.25}
\tikzset{
  femm/materials/conducting iron/.style={conductivity=2000000,
    bh={0/0,.25/25.293,.5/59.375,.75/146.191,1/400,
      1.25/1040.527,1.5/2428.125,1.75/5098.926,2/9800,
      2.25/17524.512,2.5/29546.875,3/73200}},
  biased rod/.style={femm/problem={type=axisymmetric,
      xmax=25,ymin=-25,ymax=25,mesh size=1.5,interpolation=femm},
    femm/circuit={name=drive,current=#1}},
  ac response/.style 2 args={femm/problem={operating point=rod-#1,
    frequency=#2,permeability=incremental,
    circuit currents={drive=1},circuit phases={drive=0}}},
  femm/arrow spacing=14mm}
\newcommand\RodGeometry{
  \path[femm/region={material=air,mesh size=.6}] (0,-14) rectangle (11,14);
  \path[femm/region={material=conducting iron,mesh size=\CoreMesh}]
    (0,-12) rectangle (5,12);
  \path[femm/region={material=copper,circuit=drive,turns=100,
    winding fill=.6,mesh size=.8}] (7,-8) rectangle (10,8);}
\end{codeexample}

The following definitions complete the shared setup.

% manual example: biased-harmonic/2

\begin{codeexample}[setup code]
\foreach \Case/\Bias in {low/5,high/50}{
  \setbox\BiasCalculation=\hbox{
    \begin{tikzpicture}[biased rod=\Bias,
      femm/cache={file=\jobname-bias-\Case,mode=auto},
      femm/operating point={name=rod-\Case}]
      \RodGeometry
    \end{tikzpicture}}
  \foreach \Frequency in {5,10,20,50,100,200}{
    \setbox\BiasCalculation=\hbox{
      \begin{tikzpicture}[ac response={\Case}{\Frequency},
        femm/cache={file=\jobname-\Case-\Frequency,mode=auto}]
        \femmmeasure[name=port-\Case-\Frequency,quantity=circuit,
          circuit=drive,parameter=\Frequency]
        \femmmeasure[name=loss-\Case-\Frequency,parameter=\Frequency]
      \end{tikzpicture}}}}
\end{codeexample}

The following picture uses the preceding definitions.

% manual example: biased-harmonic/3

\begin{codeexample}[width=0pt]
\foreach \Case/\Bias/\Colour in {low/5/weakbias,high/50/strongbias}{
  \begin{tikzpicture}[ac response={\Case}{200},scale=1.25,
    femm/cache={file=\jobname-\Case-200,mode=auto}]
    \clip (-3,-18) rectangle (24,22);
    \pic {femm geometry};
    \pic[femm/phase=0,femm/lines=13,femm/field style={\Colour}]
      {femm field};
    \pic[femm/currents={view=symbols,phase=0,spacing=1.5},
      femm/eddy style={color=\Colour}] {femm eddy currents};
    \node[fill=white,font=\small,align=center] at (11,19)
      {$I_0=\Bias\,\mathrm A$\\$f=200\,\mathrm{Hz}$, $\omega t=0$};
    \draw[->] (0,-15)--(12,-15) node[right,font=\scriptsize] {$r$};
    \node[font=\scriptsize,fill=white] at (2.5,-13) {Iron};
    \node[font=\scriptsize,fill=white,rotate=90] at (8.5,0) {Coil};
  \end{tikzpicture}}
\end{codeexample}

% manual example: biased-harmonic/4

\begin{codeexample}[setup code]
\def\LowPorts{port-low-5,port-low-10,port-low-20,port-low-50,port-low-100,port-low-200}
\def\HighPorts{port-high-5,port-high-10,port-high-20,port-high-50,port-high-100,port-high-200}
\def\LowLoss{loss-low-5,loss-low-10,loss-low-20,loss-low-50,loss-low-100,loss-low-200}
\def\HighLoss{loss-high-5,loss-high-10,loss-high-20,loss-high-50,loss-high-100,loss-high-200}
\end{codeexample}

The following picture uses the preceding definitions.

% manual example: biased-harmonic/5

\begin{codeexample}[width=0pt]
\begin{tikzpicture}
\begin{semilogxaxis}[width=7.5cm,height=5.7cm,grid=major,
  xlabel={Frequency (Hz)},ylabel={$\operatorname{Re}L_{\rm inc}^{*}$ (mH)},
  y filter/.expression={1000*y},scaled y ticks=false,
  legend style={font=\scriptsize,at={(.03,.03)},anchor=south west}]
  \femmresultplot[results={\LowPorts},component=incremental_inductance_real,
    plot options={weakbias,thick,mark=*}]
  \addlegendentry{$I_0=5\,\mathrm A$}
  \femmresultplot[results={\HighPorts},component=incremental_inductance_real,
    plot options={strongbias,thick,mark=square*}]
  \addlegendentry{$I_0=50\,\mathrm A$}
\end{semilogxaxis}
\end{tikzpicture}
\end{codeexample}

% manual example: biased-harmonic/6

\begin{codeexample}[width=0pt]
\begin{tikzpicture}
\begin{loglogaxis}[width=7.5cm,height=5.7cm,grid=major,
  xlabel={Frequency (Hz)},ylabel={Rod eddy-current loss (W)},
  legend style={font=\scriptsize,at={(.03,.97)},anchor=north west}]
  \femmresultplot[results={\LowLoss},component=eddy_loss,
    plot options={weakbias,thick,mark=*}]
  \addlegendentry{$I_0=5\,\mathrm A$}
  \femmresultplot[results={\HighLoss},component=eddy_loss,
    plot options={strongbias,thick,mark=square*}]
  \addlegendentry{$I_0=50\,\mathrm A$}
\end{loglogaxis}
\end{tikzpicture}
\end{codeexample}

\luafemmexampleend{biased-harmonic}

Read the line panels as instantaneous \emph{response} fields.
Their levels are chosen independently. A denser-looking panel alone does not
prove a larger flux. The dot/cross symbols indicate current direction normal
to the meridional drawing, with the convention of Section~\ref{sec:currents}.
The inductance and loss curves provide quantitative comparisons with explicit
units. The |eddy_loss| component includes the resolved rod currents and excludes
stranded-coil copper loss. Neither curve includes DC heating.
For these parameters the stronger bias lowers the real differential inductance,
but its rod loss becomes larger as frequency rises. A lower tangent permeability
also changes penetration and the current distribution: saturation alone does
not imply a monotone decrease of eddy-current loss.

\subsection{Complex inductance, impedance and power}
For a single excited port define
\[
 L_{\mathrm{inc}}^*=\frac{\widehat\lambda}{\widehat I},\qquad
 Z=R+i\omega L_{\mathrm{inc}}^*.
\]
The superscript star here labels a complex inductance; it is not a conjugation.
The linkage includes the response of the conducting body. Consequently
\[
 \operatorname{Re}Z-R=-\omega\operatorname{Im}L_{\mathrm{inc}}^*,
 \qquad \operatorname{Im}Z=\omega\operatorname{Re}L_{\mathrm{inc}}^*.
\]
With the package's sign convention a passive eddy-current load normally has
negative imaginary inductance. At low frequency the real component approaches
the static differential inductance, provided the same bias and mesh are used.
Do not replace it by $\lambda_0/I_0$, which is the static secant ratio.

% Keep the key heading with its description after the full-page figure.
\newpage
\begin{key}{/femm/problem/circuit phases=\marg{name=value list} (initially empty)}
Set selected inherited circuit phases in degrees. PGF evaluates each numerical
expression. For example, with both circuits already declared in the bias:
\begin{codeexample}[code only]
circuit phases={drive=0,sense=180/2}
\end{codeexample}
Names must exist in the captured bias or sliding-band template. This key
requires |operating point| or |sliding band|; a nonzero phase requires positive
frequency. Unlisted circuits retain their captured phases (zero for a DC bias).
An empty list changes no phase. Source amplitudes are set
separately by |circuit currents|. Both lists are ordinary styleable PGF keys.
\end{key}

For an operating-point response, |/femm/problem/frequency| is the one allowed physical
problem override. Its default remains zero (the static response); a positive
value in hertz selects the harmonic response and requires
|permeability=incremental|. |permeability=frozen| with positive frequency is
rejected: its retained DC sources must not silently become AC sources. Other
problem settings, material laws, geometry, boundaries and solver tolerances
come from the bias. Set them before capture.

The circuit result adds two components, both in henries:
\begin{itemize}
\item |incremental_inductance_real| gives the real part;
\item |incremental_inductance_imaginary| gives the imaginary part.
\end{itemize}
They are available only on
an incremental harmonic model with exactly one nonzero circuit excitation,
and only for that excited circuit. A zero-current sensing port still reports
complex linkage. With several simultaneously excited ports, use those linkages
and the declared current vector; no scalar self-inductance ratio is reported.
To build a small-signal matrix, solve one excitation at a time. The existing
|impedance_real|, |impedance_imaginary|, voltage, resistance and complex linkage
components retain their meanings and eligibility rules.
\begin{codeexample}[code only]
\femmvalue[component=incremental_inductance_real]{terminal}
\femmresultplot[results={port-low-5,port-low-10,port-low-20},
  component=incremental_inductance_real,
  plot options={blue!65!black,thick,mark=*}]
\end{codeexample}
Record the frequency with |parameter| when measuring, as in the complete
example. Profile paths retain the usual syntax; use |representation=amplitude|
for $|\widehat B_z|$, |real| or |imaginary| for phasor components, and |snapshot|
with |phase| for instantaneous fields. A harmonic component amplitude is
nonnegative; it loses the phase information needed to reconstruct a waveform.

The cycle-average response energy and terminal powers are
\[
 \overline W=\frac14\int
 \widehat{\mathbf B}^{\dagger}\boldsymbol\nu_{\mathrm{inc}}
 \widehat{\mathbf B}\,\mathrm dV,\quad
 P=\frac12\operatorname{Re}(\widehat V\overline{\widehat I}),\quad
 Q=\frac12\operatorname{Im}(\widehat V\overline{\widehat I}).
\]
The dagger denotes complex conjugate transpose. For the example's single
driven port and homogeneous Dirichlet/natural boundaries,
$Q=2\omega\overline W$, while $P$ equals the total resolved Joule loss plus
stranded-coil copper loss. With several driven ports, sum their complex powers.
Regional energy excludes any energy represented by a Robin boundary term;
include that term separately when checking reactive power for such a boundary.
Likewise, a prescribed conductor-length override can include resistance of wire
outside the model volume, whose loss is absent from the regional integral.
These are AC response quantities, excluding DC
energy and DC winding heating. Mechanical measurements separately include the
first-order bias--response cross terms described in Section~\ref{sec:mechanical-response}.

\subsection{Constitutive model, persistence and limits}
The differential tensor is evaluated from the captured $\mathbf B_0$, using
the formula in Section~\ref{sec:operating}. It is fixed throughout the complex
solve. Both real and imaginary response fields use that same tensor; the
solver never evaluates a new B--H tangent from the AC amplitude. In axisymmetry
it evaluates the local bias curl and tensor at each volume quadrature point,
including the $A_\phi/r$ term. The stiffness is the static differential
operator; the conductive term is $i\omega\sigma A$. In this formulation
$\widehat J_{\rm eddy}=-i\omega\sigma\widehat A$.

Permanent magnets may establish the DC bias. Their coercive source is removed
from the incremental excitation, while their supported recoil law remains.
A real tangent does not model hysteresis, demagnetisation history, complex
permeability or frequency-dependent material losses. Nor is this a nonlinear
large-amplitude harmonic or transient solver.

The captured operating-point interface retains prescribed series-current
windings, real positive material tensors, supported ordinary boundaries, and
no floating-conductor constraints or voltage/parallel drives in these responses,
harmonic homogenised laminations or circular open boundaries. Unbiased harmonic circuits support the
additional constraints in Section~\ref{sec:ac-circuits}. Here planar solid conductors use zero imposed
longitudinal electric-potential gradient; axisymmetric currents close naturally.
Supported imported DC answers can supply a bias (Section~\ref{sec:interchange});
planar response-model FEM/ANS exchange is described in
Section~\ref{sec:response-exchange}.

Use a distinct cache file for each bias/frequency/excitation. Native cache
format 24 includes the exact bias mesh and potential, mode, frequency and
source amplitudes/phases. Presentation-only phase changes reuse it. The named
bias is still captured in each TeX run, optionally from its own restored DC
cache. There is no new persistent format for the captured name itself.

\subsection{Lua entry point and validation}
The same numerical API accepts frequency and phase overrides:
\begin{codeexample}[code only]
local response = require("luafemm-operating").new(bias, {
  mode = "incremental", frequency = 200,
  currents = {drive = 1}, phases = {drive = 90},
  cache = {file = "biased-200", mode = "auto"}
})
require("luafemm").solve(response)
local port = require("luafemm-harmonic").circuit(response, "drive")
\end{codeexample}
The bias must be a native converged unchanged DC model or its owned capture.
Frequency is a finite nonnegative number in hertz; phases are finite numbers
in degrees. Frequency and phases default to zero and an empty table.

Regression checks cover an analytic biased conducting slab with refinement,
low-frequency limits, phase covariance, permanent-magnet source removal,
power/energy balances and immutable bias/cache data. An optional xfemm oracle
solves independent linear AC problems with analytically derived tangent tensors
at a manufactured uniform bias. It checks the differential constitutive and
complex diffusion operators without claiming support for xfemm's previous-field
file workflow. The generic format matrix exercises the same phasor caches in
plain LuaTeX, LuaLaTeX and ConTeXt MkIV. Quantitative results and tolerances are
recorded in \nolinkurl{doc/validation.md}.
